% !TeX root = ../drafts/04-committing-the-witness.tex
\section{Committing the witness}\label{sec:rs}

\subsection{Jagged commitment layout}\label{sec:stacking}

No ordinary trace column is committed on its own. We use the jagged PCS construction of Hemo et al.~\cite{Jagged26}, written out in Annex~\ref{annex:jagged}. Compatible equal-height columns are partitioned into power-of-two-width blocks (the ``fancy jagged'' layout of that work): compatibility means that the columns occur in exactly the same opening groups. For a block of width $B=2^b$ and real height $h$, its columns are stored row-major,
\[
  (P_0[0],\ldots,P_{B-1}[0],P_0[1],\ldots,P_{B-1}[h-1]).
\]
The blocks themselves are tightly concatenated and the final vector is zero-padded only once, to $2^M$. A non-power-of-two compatible group is decomposed into power-of-two blocks, so this table packing adds no dummy cells. Opcode columns use their announced row counts. The memory image and memory-finalize count use the announced used-memory prefix, with public suffixes $0$ and $1$; the bytecode-finalize count uses the program-bound real instruction prefix, with public suffix $1$. Thus $q_\flock$ is the only full power-of-two committed column. If $A$ is the sum of the physical block lengths, then
\[
  M=\max\!\left(15,\lceil\log_2\max(A,1)\rceil,
                    \max_{\text{blocks}}(k+b)\right),
\]
where $k$ is a block's logical row dimension and $b$ its selector width. The last term ensures that every logical point embeds in the dense cube even when its physical prefix is short.

If such a block begins at dense offset $t$, its column $c$ and row $z$ correspond to the Boolean dense index
\[
  w=t+c+Bz \quad\text{with}\quad w<t+Bh.
\]
The low $b$ coordinates of the logical point are fixed to the bits of $c$ and the remaining coordinates are the row point. Let $J_{t,Bh}((c,z),w)$ be the multilinear extension of this predicate. It is the width-four read-once branching program of Annex~\ref{annex:jagged}: its state records the carry of $t+(c,z)$ and the most-significant comparison seen between $w$ and $t+Bh$, while reading bits least-significant first. The real part of a column evaluation is consequently the inner-product claim
\[
  \sum_{w\in\cube M} J_{t,Bh}((c,\zeta),w)\,\widetilde q(w).
\]
The arithmetization still regards a column as extended to its power-of-two logical height with a fixed public pad value $p_i$. The verifier adds the omitted suffix itself: if $H_i(\zeta)$ is the MLE of $[z<h_i]$, the logical claim is the real contribution plus $p_i(1-H_i(\zeta))$. (Addition and subtraction coincide in characteristic two.)

The packed flock witness $q_\flock$ is the sole aligned exception. It occupies the offset-zero prefix so its ring-switch and fixed-slot claims retain their direct subcube formulas; all ordinary trace columns follow it in the tight jagged layout.

\paragraph{Batching.} Every claim the protocol raises against a column is such a weighted sum $\sum_w W_j(w)\,\widetilde q(w)=c_j$. The verifier folds the $J$ of them with powers of one random $\lambda$ into one weight $W_\lambda=\sum_j\lambda^j W_j$ and target $C_\lambda=\sum_j\lambda^j c_j$, and the inner-product PCS (Annex~\ref{annex:pcs}) proves
\[
  \sum_{w\in\cube{M}} W_\lambda(w)\,\widetilde q(w)=C_\lambda.
\]
Claims covering a complete width-$B$ block receive consecutive powers in column-selector order. Their sum is one jagged indicator whose selector-coordinate weights are the unnormalised pairs $(1,\lambda^{2^r})$, scaled by the first power assigned to the block. Thus the verifier runs one branching program per claimed block rather than per physical column, with no field inversions. Unlike the generic adapter in~\cite{Jagged26}, our underlying PCS already accepts a verifier-evaluable inner-product weight, so we evaluate these batched indicators directly and introduce no separate jagged reduction sumcheck.

\subsection{Committing to booleans via Ring-Switching}\label{sec:ringswitch}

Flock's BLAKE3 argument (\S\ref{sec:tab-blake3}) requires to commit a boolean multilinear (i.e. taking values in $\Ftwo$ on the boolean hypercube) in $\nflock + 6$ variables. Ring-Switching~\cite{DP26} enables us to transform a commitment in $\K$ into a commitment in $\Ftwo$, see Annex~\ref{annex:rs}. As a result, the Flock witness is committed as a $\K$-valued multilinear in $\nflock$ variables (using stacking \ref{sec:stacking}).

